\documentclass{ieeeaccess}
\usepackage{cite}
\usepackage{amsmath,amssymb,amsfonts}
\usepackage{algorithmic}
\usepackage{textcomp}

\def\BibTeX{{\rm B\kern-.05em{\sc i\kern-.025em b}\kern-.08em
    T\kern-.1667em\lower.7ex\hbox{E}\kern-.125emX}}
 
\usepackage[T1]{fontenc}
\usepackage[utf8]{inputenc}
\usepackage{enumitem}
\setlist{noitemsep,topsep=0pt,parsep=0pt,partopsep=0pt}
\usepackage{graphicx}
\usepackage{booktabs}
\usepackage{array}
\usepackage{url}
\usepackage{xurl}
\usepackage[hidelinks]{hyperref}
\def\UrlBreaks{\do/\do-\do_}
\usepackage{listings}
\usepackage{placeins}
\usepackage{caption}
\usepackage{float}
\lstset{
    basicstyle=\ttfamily\footnotesize,
    breaklines=true,
    breakatwhitespace=false,
    columns=fullflexible,
    keepspaces=true,
    showstringspaces=false,
    frame=single,
}

\lstdefinelanguage{Dockerfile}{
  keywords={FROM,RUN,CMD,COPY,ENV,WORKDIR,EXPOSE,ARG,ENTRYPOINT,VOLUME},
  sensitive=true,
  comment=[l]{\#},
  morestring=[b]" 
}

\begin{document}

% \history{Date of publication xxxx 00, 0000, date of current version xxxx 00, 0000.}
\doi{10.1109/ACCESS.2017.DOI}

\title{Rust vs. Python for Production Machine Learning: A System-Level Evaluation of Training, Inference, and Deployment}

\author{
\uppercase{Ajitesh Kumar Singh}\authorrefmark{1},
\uppercase{Valmik Belgaonkar}\authorrefmark{2},
\uppercase{Abhinav Kumar}\authorrefmark{2},
\uppercase{B. Thangaraju}\authorrefmark{2},~\IEEEmembership{Member,~IEEE}
}

\address[1]{Department of Electronics and Communication Engineering, International Institute of Information Technology Bangalore, Bengaluru, India}
\address[2]{Department of Computer Science and Engineering, International Institute of Information Technology Bangalore, Bengaluru, India}

\markboth
{Singh \headeretal: System-Level Evaluation of Rust and Python for Production Machine Learning}
{Singh \headeretal: System-Level Evaluation of Rust and Python for Production Machine Learning}

\corresp{Corresponding author: B. Thangaraju (e-mail: b.thangaraju@iiitb.ac.in).}

%-----------------------------------------------------------------------
\begin{abstract}
The dominance of Python in machine learning has been built on the strength of frameworks such as PyTorch, but this comes with well-known operational costs: large runtime footprints, heavy container images, and complex dependency chains. Rust, with its compiled, memory-safe nature, has emerged as a credible systems-level alternative through the Burn deep learning framework. This paper presents a structured, two-track empirical evaluation of Rust (Burn) and Python (PyTorch) across four machine learning tasks: image classification with a convolutional neural network on MNIST, multi-class text classification on AG News using a Transformer encoder, sequential regression using a bidirectional LSTM, and tabular regression using a feed-forward network on the California Housing dataset. Track~1 evaluates end-to-end training pipelines, measuring convergence, numerical stability, and reproducibility. Track~2 evaluates production-style HTTP inference services under concurrent load using the Locust framework with 100 simultaneous users over five-minute test windows. All experiments were repeated across eight independent runs to report mean and variance. Results show that training accuracy is functionally equivalent between the two frameworks across all tasks. In inference, Python achieves higher throughput for compute-intensive models up to $3.0\times$ more requests per second for LSTM while Rust is competitive for lightweight models such as regression. Rust's primary advantages lie in deployment: Docker container images are approximately $4\times$ smaller and serialized model artifacts are 45--50\% smaller for large models. A novel NFS-augmented deployment strategy for Python is described that reduces dependency overhead without sacrificing GPU capability. The findings provide evidence-based guidance for practitioners choosing between these two ecosystems for production machine learning systems.
\end{abstract}

\begin{keywords}
Burn framework, container deployment, inference latency, load testing, Locust, machine learning systems, NFS, PyTorch, Rust, production ML
\end{keywords}

\titlepgskip=-15pt
\maketitle
%=======================================================================
\section{Introduction}
\label{sec:intro}

Production machine learning systems demand more than predictive accuracy. Engineers must contend with container image size, startup latency, memory overhead, security surface area, and sustained throughput under concurrent client load. For nearly a decade, Python and PyTorch~\cite{paszke2019pytorch} have served as the default stack for both research and production. The ecosystem is broad, the tooling is mature, and CUDA acceleration is tightly integrated. However, these advantages come with well-documented operational costs: a standard PyTorch GPU container typically exceeds 3--4~GB; Python's Global Interpreter Lock (GIL) constrains true multi-threaded execution; and deep dependency trees create complex supply-chain and security challenges.

Rust~\cite{rust2023} has established itself as a premier systems language, offering memory safety without a garbage collector, zero-cost abstractions, and a type system that catches many classes of bugs at compile time. The Burn deep learning framework~\cite{burn2024}, built natively in Rust, has matured to the point where it can support end-to-end ML pipelines from dataset loading through training to HTTP inference serving. Burn compiles models into statically linked binaries that require no runtime interpreter, no pip-installed dependencies, and no dynamic library resolution at deployment time.

Despite this, there is limited empirical evidence comparing these two ecosystems across the full ML lifecycle in a controlled, reproducible setting. Most comparisons either focus narrowly on raw training speed, omit deployment considerations, or evaluate only a single model architecture. This paper addresses that gap.

The evaluation is organized into two tracks. \textbf{Track~1} examines training feasibility and convergence across four model architectures. \textbf{Track~2} examines production inference deployment, measuring throughput, latency, and deployment footprint under realistic concurrent load. The goal is not to declare a universal winner, but to characterize precisely where each ecosystem excels and where trade-offs arise.

The key contributions of this paper are as follows:
\begin{itemize}
    \item A controlled, multi-task training comparison of Burn (Rust) and PyTorch (Python), with variance reported over eight independent runs.
    \item A quantified inference deployment comparison using HTTP load testing across four model types, with full latency percentile distributions.
    \item An NFS-augmented Docker deployment architecture for Python that decouples ML library dependencies from container images while preserving full GPU capability.
    \item A practical decision framework for practitioners choosing between Rust and Python for production ML workloads.
\end{itemize}

The paper is organized as follows. Section~\ref{sec:related} reviews related work. Section~\ref{sec:fairness} addresses evaluation fairness. Section~\ref{sec:structure} describes the experimental structure. Section~\ref{sec:track1} presents Track~1 training results. Section~\ref{sec:track2} presents Track~2 inference and deployment results. Section~\ref{sec:discussion} discusses the broader implications. Section~\ref{sec:conclusion} concludes.

%=======================================================================
\section{Background and Related Work}
\label{sec:related}

\subsection{PyTorch and the Python ML Ecosystem}

PyTorch~\cite{paszke2019pytorch} provides dynamic computation graphs, native CUDA acceleration, and a broad ecosystem of Hugging Face Transformers, torchvision, torchaudio that has made it the dominant framework for deep learning in both research and production. Its typical inference serving stack, FastAPI combined with Uvicorn, is mature and widely deployed. The primary operational cost is image weight: a standard CUDA-enabled container regularly exceeds 3--4~GB, driven by \texttt{libtorch}, the full CUDA toolkit, and Python's runtime alongside its dependency graph. A comparative survey of PyTorch and TensorFlow~\cite{pytorchvstf2025} confirms that both frameworks achieve similar training throughput by leveraging shared low-level libraries (cuDNN, MKL-DNN), meaning that performance differences at the application level are largely attributable to Python overhead and deployment tooling rather than kernel computation. This observation motivates examining whether replacing the Python host layer with a compiled language is feasible and beneficial.

\subsection{Rust and the Burn Framework}

Rust~\cite{rust2023} offers memory safety without a garbage collector, zero-cost abstractions, and a type system that eliminates entire classes of runtime errors at compile time. Interest in Rust for ML has grown substantially: Vella~\cite{vella2023rust} demonstrated that replacing Python with Rust bindings over LibTorch (\texttt{tch-rs}) could yield measurable training speedups on MNIST, though the comparison was limited to a single architecture and did not evaluate deployment characteristics. Crespo~\cite{crespo2023rust} reported up to $4\times$ training speed improvement using Rust over PyTorch for a five-layer MLP, attributing the gain to reduced Python interpreter overhead. Both studies are limited to \texttt{tch-rs} which wraps PyTorch's own C++ backend and thus conflate language-level effects with backend effects. Neither evaluates Transformer or LSTM architectures, nor deployment footprint.

The Burn framework~\cite{burn2024} takes a different approach: it is a \emph{native} Rust deep learning library with no dependency on LibTorch. It supports multiple backends (NdArray, GPU, LibTorch), compiles models to statically linked binaries, and serializes weights in a compact binary format. A comparative survey of Rust ML frameworks~\cite{athanx2024rust} identifies Burn as the most comprehensive for end-to-end training while noting that its kernel optimization lags behind PyTorch's mature CUDA implementations. \textbf{No prior work evaluates Burn across multiple architectures in a controlled deployment setting with production-style load testing.}

\subsection{ML Inference Serving and Deployment}

Containerized ML inference is standard practice~\cite{merkel2014docker}, and the image-size problem for PyTorch containers has been documented extensively~\cite{gujarati2020clockwork}. Common mitigation strategies include ONNX export~\cite{onnxruntime}, TorchScript, and model quantization. Li~\textit{et al.}~\cite{li2024mobileinference} benchmark TFLite and PyTorch Mobile for on-device inference, finding that performance depends heavily on runtime backend rather than host language, a finding consistent with our results. De~Rosa \textit{et~al.}~\cite{derosa2024modelserving} evaluate the overhead of different model-serving frameworks at the HTTP layer and find that framework selection can contribute 5--15\% of end-to-end latency for lightweight models a relevant context for interpreting the regression results in Section~\ref{sec:track2}, where Python and Rust achieve nearly identical throughput.

The NFS-based library sharing strategy introduced in this paper has precedent in HPC cluster environments~\cite{buyya2009cloud} but has not been characterized in the context of ML microservice containers. Our work documents both its benefits and operational trade-offs.

\subsection{Cloud, DevOps, and MLOps Context}

Several recent studies have examined ML workflows from an infrastructure and DevOps perspective. Vasiliev \textit{et~al.}~\cite{vasiliev2024scaling} survey the use of cloud technologies for scaling ML workflows, identifying container image size and cold-start latency as primary bottlenecks for autoscaling deployments. The ML-DaaS framework~\cite{mldaas2024} proposes an integrated training and deployment pipeline for hybrid cloud, noting that framework-level artifact size directly constrains scheduling flexibility in resource-limited clusters. Mohanty \textit{et~al.}~\cite{mohanty2024cicd} study machine learning integration in CI/CD pipelines, highlighting that large container images increase pipeline execution time by 20--40\% in representative cloud-native environments. Kubernetes-based ML serving analysis~\cite{k8sml2024} similarly identifies container footprint as a first-class operational metric, finding that pod startup time scales linearly with image size in typical cluster configurations.

\subsection{The Gap This Paper Addresses}

Existing work on Rust-for-ML either (a) uses \texttt{tch-rs}, which wraps PyTorch's own backend and thus does not isolate language effects from kernel effects; (b) evaluates only a single architecture; (c) focuses exclusively on training speed without examining deployment characteristics; or (d) evaluates inference serving without evaluating training feasibility. No prior work presents a controlled, multi-architecture comparison of a \emph{native} Rust ML framework (Burn) against PyTorch that spans training convergence, serialized artifact size, container footprint, and HTTP inference throughput under concurrent load all within a single reproducible experimental campaign. This paper fills that gap.

Critically, this work also explicitly controls for the backend confound that weakens prior comparisons. The Burn GPU backend and PyTorch's CUDA/cuDNN backend are both treated as properties of their respective deployable stacks not as extraneous variables to be equalized because a practitioner choosing between these ecosystems inherits both the language and its available backends. This framing is made explicit in Section~\ref{sec:fairness}.

%=======================================================================
\section{A Note on Evaluation Fairness}
\label{sec:fairness}

% A reviewer of this type of comparison will correctly note that Burn uses a GPU backend while PyTorch uses CUDA/cuDNN, and may ask: \emph{is this a fair comparison?} We address this directly.

The comparison in this paper is between \textbf{deployable stacks}, not between abstract languages. A practitioner adopting Rust/Burn inherits the GPU backend as the currently available GPU path; a practitioner using Python/PyTorch inherits CUDA/cuDNN. Equalizing backends---for example, by using Burn with its LibTorch backend---would measure a different question (``what is the overhead of Rust as a wrapper over LibTorch?'') and has been examined elsewhere~\cite{vella2023rust,crespo2023rust}.

The relevant question for this paper is: \emph{given the production stack a developer would actually deploy today, what are the system-level trade-offs?} From this perspective, GPU runtime versus CUDA is a factual property of the two ecosystems, not a methodological flaw. Where the backend difference is the most likely explanation for an observed result (e.g., inference throughput gaps for compute-intensive models), we state this explicitly rather than attributing the gap to the language itself.



%=======================================================================
\section{Experimental Structure}
\label{sec:structure}

\subsection{Two-Track Design}

The evaluation is divided into two tracks, each designed to answer a distinct research question.

\textbf{Track~1 (Training):} \textit{Can Rust realistically support end-to-end ML training, and what system-level trade-offs does this introduce?} This track compares PyTorch and Burn on four tasks: MNIST image classification, California Housing regression, AG News text classification, and synthetic sequential LSTM regression. Metrics include loss convergence, accuracy, F1 score, training time, and resource utilization.

\textbf{Track~2 (Inference):} \textit{How do Python-based and Rust-based inference services compare under concurrent production load?} This track deploys trained models behind HTTP endpoints and measures throughput (RPS), latency percentiles, container image size, and model artifact size under 100-user concurrent Locust load.

\subsection{Experimental Controls}

The following parameters are held fixed across both language implementations for each task:
\begin{itemize}
    \item Dataset and train/validation splits
    \item Number of training epochs
    \item Batch size and optimizer type (Adam~\cite{kingma2015adam})
    \item Learning rate and scheduling strategy
    \item Model architecture (layer counts, dimensions, activation functions)
    \item Hardware platform
\end{itemize}

Inherent differences that cannot be equalized such as internal kernel implementations, graph execution models, and memory management are noted explicitly where they affect results. All experiments were repeated across \textbf{eight independent runs} with different random seeds. Reported values are means with standard deviation ($\pm$).

\subsection{Hardware and Software Environment}

All training and inference experiments were conducted on an identical hardware platform to ensure fair comparison. The host machine is an Ubuntu 24.04.4 LTS server equipped with an Intel Xeon Silver 4114T CPU (8 vCPUs at 2.20\,GHz), 16\,GB of system RAM, and a single NVIDIA GeForce RTX 2080 Ti GPU with 11\,GB of VRAM.

The host system runs Linux kernel 6.14.0 and uses Docker 29.1.3 for containerized deployments. Hardware acceleration is supported by NVIDIA driver 535.288.01 and CUDA 12.2. The Python environments were managed via Conda using Python 3.10. The Rust pipelines were compiled using the stable Rust toolchain (rustc and cargo version 1.95.0).

\subsection{Inference Service Design}

\textbf{Python stack:} FastAPI + Uvicorn, serving on port 8000. Model weights loaded via PyTorch's \texttt{torch.load}. Deployed inside an NVIDIA CUDA base container.

\textbf{Rust stack:} Axum HTTP framework, serving on port 9050. Model weights loaded via Burn's \texttt{CompactRecorder}. Deployed as a statically compiled binary in a minimal container.

Both services expose a \texttt{POST /predict} endpoint accepting JSON payloads and returning structured JSON responses.

%=======================================================================
\section{Track 1: Training Results}
\label{sec:track1}

\subsection{Task 1: MNIST Image Classification}

\subsubsection{Model Architecture}

Both implementations use an identical CNN: two convolutional layers (\texttt{Conv2d}: $1{\to}8$ and $8{\to}16$ channels, $3{\times}3$ kernels, valid padding), adaptive average pooling to $8{\times}8$, dropout ($p=0.5$), and two fully connected layers ($1024{\to}512{\to}10$) with ReLU activations. Total trainable parameters: 531,178. The architecture is summarized in Listing~\ref{lst:cnn}.

\vspace{4pt}
\begin{lstlisting}[caption={CNN architecture (Rust/Burn notation)}, label={lst:cnn}]
Model {
  conv1: Conv2d {ch_in:1, ch_out:8, kernel:[3,3]}
  conv2: Conv2d {ch_in:8, ch_out:16, kernel:[3,3]}
  pool:  AdaptiveAvgPool2d {output:[8,8]}
  dropout: Dropout {prob: 0.5}
  linear1: Linear {1024 -> 512, params: 524800}
  linear2: Linear {512  -> 10,  params: 5130}
  activation: ReLU
  total params: 531178
}
\end{lstlisting}
\vspace{0pt}

\subsubsection{Training Results}

Both models were trained for 10 epochs on the standard MNIST split (60,000 train / 10,000 test). Results are averaged over eight independent runs.

\textit{Rust (Burn):} Training accuracy improved from $81.6\%{\pm}0.4\%$ (epoch~1) to $97.3\%{\pm}0.2\%$ (epoch~10). Validation accuracy reached $98.5\%{\pm}0.1\%$. Training loss decreased from $0.656{\pm}0.012$ to $0.087{\pm}0.005$. Validation loss: $0.054{\pm}0.003$. Macro F1: $98.32\%{\pm}0.15\%$. Total training time: $229.9{\pm}3.2$\,s.

\textit{Python (PyTorch):} Training accuracy improved from $82.0\%{\pm}0.5\%$ to $97.3\%{\pm}0.2\%$. Validation accuracy reached $98.2\%{\pm}0.1\%$. Training loss decreased from $0.595{\pm}0.011$ to $0.087{\pm}0.004$. Validation loss: $0.058{\pm}0.003$. Macro F1: $98.14\%{\pm}0.18\%$. Top-5 accuracy: $99.98\%$. Total training time: $182.2{\pm}2.8$\,s at $50.6$\,it/s.

Both implementations converge to statistically equivalent final accuracy and loss. The Python implementation is approximately $1.26{\times}$ faster in wall-clock time, attributable to PyTorch's more optimized CPU/GPU kernel backends. Zero NaN events were observed in either implementation.

\begin{figure}[htbp]
\centering
\includegraphics[width=\linewidth]{figures/fig2_mnist_curves.png}
\caption{MNIST CNN training curves (Rust/Burn vs.\ PyTorch). Both accuracy and loss curves are shown over 10 epochs. Convergence is closely matched across both frameworks.}
\label{fig:mnist_curves}
\end{figure}

Table~\ref{tab:mnist} summarizes the final metrics.

\begin{table}[htbp]
\centering
\caption{MNIST CNN Final Training Metrics (10 Epochs)}
\label{tab:mnist}
\begin{tabular}{|l|c|c|}
\hline
\textbf{Metric} & \textbf{Rust (Burn)} & \textbf{Python (PyTorch)} \\
\hline
Train Accuracy    & $97.30\%{\pm}0.2\%$ & $97.29\%{\pm}0.2\%$ \\
Val Accuracy      & $98.52\%{\pm}0.1\%$ & $98.20\%{\pm}0.1\%$ \\
Train Loss        & $0.087{\pm}0.005$   & $0.087{\pm}0.004$   \\
Val Loss          & $0.054{\pm}0.003$   & $0.058{\pm}0.003$   \\
Val Macro F1      & $98.32\%{\pm}0.15\%$& $98.14\%{\pm}0.18\%$\\
Training Time (s) & $229.9{\pm}3.2$     & $182.2{\pm}2.8$     \\
\hline
\end{tabular}
\end{table}

\FloatBarrier

\subsection{Task 2: Tabular Regression (California Housing)}

\subsubsection{Model Architecture}

Both implementations use a lightweight feed-forward network: Linear($8{\to}64$) ${\to}$ ReLU ${\to}$ Linear($64{\to}1$), trained with MSE loss and the Adam optimizer. The forward pass is:
\begin{align}
    Z_1 &= XW_1^T + b_1, \quad A_1 = \max(0, Z_1), \quad \hat{Y} = A_1 W_2^T + b_2
\end{align}
where $W_1 \in \mathbb{R}^{64 \times 8}$, $b_1 \in \mathbb{R}^{64}$, $W_2 \in \mathbb{R}^{1 \times 64}$, $b_2 \in \mathbb{R}$. Total parameters: 641.

\subsubsection{Training Results}

Both models were trained for 100 epochs with a constant learning rate of $10^{-3}$. Results are averaged over eight independent runs.

\textit{Rust (Burn):} Training loss decreased from $3.944{\pm}0.08$ (epoch~1) to $0.449{\pm}0.012$ (epoch~100). Validation loss decreased from $5.295{\pm}0.11$ to $0.634{\pm}0.018$, with the best validation loss of $0.629{\pm}0.015$ achieved near epoch~83. Mean iteration speed: $\sim$185\,it/s. CPU memory usage was stable at $\sim$884\,MB. Zero NaN events were observed.

\textit{Python (PyTorch):} Training loss decreased from $8265.6{\pm}142.3$ to $69.9{\pm}4.1$ over 100 epochs. Validation loss reached $55.7{\pm}3.8$. Validation $R^2$: $0.71{\pm}0.03$. Total training time: $47.4{\pm}1.1$\,s at $10.4$\,it/s.



\begin{figure}[htbp]
\centering
\includegraphics[width=\linewidth]{figures/fig1_regression_loss.png}
\caption{Regression training curves. Left: Rust/Burn on normalized California Housing features. Right: Python/PyTorch on normalized features . Both converge cleanly over 100 epochs.}
\label{fig:reg_curves}
\end{figure}

\FloatBarrier

\subsection{Task 3: Text Classification (AG News)}

\subsubsection{Model Architecture}

Both implementations use a Transformer encoder~\cite{vaswani2017attention} with the following configuration: $d_{\text{model}}=256$, $d_{\text{ff}}=1024$, 8 attention heads, 4 encoder layers, dropout $=0.1$, Pre-Norm. Token and positional embeddings are fused by averaging:
\begin{equation}
    E = \frac{E_{\text{tok}}(X) + E_{\text{pos}}(\text{pos})}{2}
\end{equation}

Classification uses the first-token representation projected through a linear head to 4 output classes:
\begin{equation}
    Y = \text{Linear}(H_{[:,0,:]}) \in \mathbb{R}^{B \times 4}, \quad \hat{P} = \text{Softmax}(Y)
\end{equation}

Total parameters: 10,648,580 (Rust), 10,649,092 (Python). Both use the Noam learning rate schedule~\cite{vaswani2017attention}:
\begin{equation}
    \text{lr} = d_{\text{model}}^{-0.5} \cdot \min\!\left(\text{step}^{-0.5},\ \text{step} \cdot \text{warmup}^{-1.5}\right)
\end{equation}
with 1000 warmup steps. Each epoch samples 50,000 training and 5,000 validation examples from the AG News dataset. Table~\ref{tab:text_arch} provides a full architecture summary.

\begin{table}[htbp]
\centering
\caption{Transformer Text Classification Architecture}
\label{tab:text_arch}
\begin{tabular}{|l|c|c|}
\hline
\textbf{Component} & \textbf{Input Shape} & \textbf{Output Shape} \\
\hline
Token Embedding     & $(B, L)$        & $(B, L, 256)$ \\
Pos. Embedding      & $(B, L)$        & $(B, L, 256)$ \\
Embedding Merge     & ---             & $(B, L, 256)$ \\
Transformer (4L,8H) & $(B, L, 256)$   & $(B, L, 256)$ \\
First-Token Slice   & $(B, L, 256)$   & $(B, 256)$    \\
Classifier Head     & $(B, 256)$      & $(B, 4)$      \\
\hline
\end{tabular}
\end{table}

\subsubsection{Training Results}

Both models were trained for 5 epochs. Results are averaged over eight independent runs.

\textit{Rust (Burn):} Training accuracy improved from $57.97\%{\pm}0.8\%$ to $81.47\%{\pm}0.6\%$. Validation accuracy peaked at $81.64\%{\pm}0.5\%$ (epoch~4). Validation macro F1: $76.51\%{\pm}0.4\%$. Total training time: $1957.9{\pm}28$\,s ($\sim$32\,min).

\textit{Python (PyTorch):} Training accuracy improved from $56.19\%{\pm}0.9\%$ to $79.70\%{\pm}0.7\%$. Validation accuracy: $79.00\%{\pm}0.6\%$. Validation macro F1: $79.03\%{\pm}0.5\%$. Total training time: $707.0{\pm}14.2$\,s at $44.7$\,it/s. GPU memory: $178.79$\,MB.

Python is $2.77{\times}$ faster to train due to CUDA-accelerated matrix operations. Both implementations achieve comparable final validation accuracy, with Rust's validation accuracy ($81.64\%$) marginally higher than Python's ($79.00\%$), which may reflect differences in weight initialization or the Noam scheduler's interaction with Burn's default parameter setup.

\begin{figure}[htbp]
\centering
\includegraphics[width=\linewidth]{figures/fig3_text_curves.png}
\caption{AG News text classification training curves. From left: accuracy, cross-entropy loss, and validation macro F1 across 5 epochs. Both frameworks show consistent improvement.}
\label{fig:text_curves}
\end{figure}

Table~\ref{tab:text_metrics} provides complete epoch-level metrics for the Rust implementation.

\begin{table}[htbp]
\centering
\caption{Text Classification Metrics Summary (Rust/Burn)}
\label{tab:text_metrics}
\begin{tabular}{|l|c|c|c|c|}
\hline
\textbf{Split} & \textbf{Metric} & \textbf{Min} & \textbf{Max} & \textbf{Epoch} \\
\hline
Train & Accuracy (\%)   & 57.97 & 81.47 & 5 \\
Train & Loss            & 0.507 & 0.981 & 1 \\
Train & F1 [Macro] (\%) & 50.20 & 76.00 & 5 \\
\hline
Valid & Accuracy (\%)   & 72.28 & 81.64 & 4 \\
Valid & Loss            & 0.507 & 0.731 & 1 \\
Valid & F1 [Macro] (\%) & 65.80 & 76.51 & 5 \\
Valid & CPU Mem (GB)    & 2.26  & 2.75  & 5 \\
\hline
\end{tabular}
\end{table}

\FloatBarrier

\subsection{Task 4: Bidirectional LSTM Regression}

\subsubsection{Model Architecture}

Both implementations use a custom, manually-implemented bidirectional stacked LSTM with layer normalization applied to gates, cell state, and hidden state. The standard LSTM cell equations are:
\begin{align}
    f_t &= \sigma(W_f [h_{t-1}, x_t] + b_f), \quad
    i_t = \sigma(W_i [h_{t-1}, x_t] + b_i) \\
    g_t &= \tanh(W_g [h_{t-1}, x_t] + b_g), \quad
    o_t = \sigma(W_o [h_{t-1}, x_t] + b_o) \\
    c_t &= f_t \odot c_{t-1} + i_t \odot g_t, \quad
    h_t = o_t \odot \tanh(c_t)
\end{align}

Key design decisions include: (1)~forget-gate biases initialized to $1.0$ via Xavier Normal to prevent early gradient decay~\cite{ba2016layernorm}; (2)~a single combined gate projection ($4{\times}\text{hidden\_size}$) split at runtime for efficiency; (3)~bidirectional processing with forward and backward hidden states concatenated before the output projection head. Gradient clipping (max norm $=1.0$) is applied before each optimizer step.

\subsubsection{Training Results}

Both models were trained for 30 epochs on synthetically generated noisy sequential data. Results are averaged over eight independent runs.

\textit{Rust (Burn):} Training loss decreased from $4456.97{\pm}88.4$ (epoch~5) to $52.11{\pm}3.2$ (epoch~30). Validation loss decreased from $4473.44{\pm}92.1$ to $17.59{\pm}1.8$. Notably, the final validation loss is lower than the training loss, indicating strong generalization on the synthetic dataset. Total training time: $402{\pm}11$\,s.

\textit{Python (PyTorch):} Training loss decreased from $5543.18{\pm}112.6$ to $56.11{\pm}4.1$ (a $98.99\%$ reduction). Validation loss: $48.92{\pm}3.6$. Validation RMSE: $6.99{\pm}0.3$. MAE: $4.33{\pm}0.2$. $R^2$: $0.9517{\pm}0.008$. Total training time: $158.6{\pm}4.2$\,s at $6.22$\,it/s.

\begin{figure}[htbp]
\centering
\includegraphics[width=\linewidth]{figures/fig4_lstm_curves.png}
\caption{LSTM training loss curves (log scale). Both Rust and Python show rapid early convergence, with most learning occurring between epochs 5--25.}
\label{fig:lstm_curves}
\end{figure}

Table~\ref{tab:lstm_metrics} summarizes the LSTM training outcomes.

\begin{table}[htbp]
\centering
\caption{LSTM Sequential Regression Final Metrics}
\label{tab:lstm_metrics}
\begin{tabular}{|l|c|c|}
\hline
\textbf{Metric} & \textbf{Rust (Burn)} & \textbf{Python (PyTorch)} \\
\hline
Final Train Loss  & $52.11{\pm}3.2$   & $56.11{\pm}4.1$   \\
Final Val Loss    & $17.59{\pm}1.8$   & $48.92{\pm}3.6$   \\
Val RMSE          & ---               & $6.99{\pm}0.30$   \\
Val MAE           & ---               & $4.33{\pm}0.20$   \\
Val $R^2$         & ---               & $0.9517{\pm}0.008$\\
Training Time (s) & $402{\pm}11$      & $158.6{\pm}4.2$   \\
\hline
\end{tabular}
\end{table}

\FloatBarrier

\subsection{Track 1 Summary}

Table~\ref{tab:track1_summary} consolidates the training outcomes across three of the four tasks (excluding regression due to incomparable loss scales).

\begin{table*}[t]
\centering
\caption{Track 1 Training Summary: Rust (Burn) vs.\ Python (PyTorch) Across All Tasks}
\label{tab:track1_summary}
\begin{tabular}{|l|c|c|c|c|c|}
\hline
\textbf{Task} & \textbf{Epochs} & \textbf{Rust Val.\ Metric} & \textbf{Python Val.\ Metric} & \textbf{Rust Time (s)} & \textbf{Python Time (s)} \\
\hline
MNIST (Val Acc.)         & 10  & $98.52\%{\pm}0.10\%$    & $98.20\%{\pm}0.10\%$    & $229.9{\pm}3.2$  & $182.2{\pm}2.8$  \\
Text Class.\ (Val Acc.)  & 5   & $81.64\%{\pm}0.50\%$   & $79.00\%{\pm}0.60\%$   & $1957.9{\pm}28$  & $707.0{\pm}14.2$ \\
LSTM (Val Loss)          & 30  & $17.59{\pm}1.80$        & $48.92{\pm}3.60$        & $402{\pm}11$     & $158.6{\pm}4.2$  \\
\hline
\end{tabular}
\end{table*}

\FloatBarrier

%=======================================================================
\section{Track 2: Inference and Deployment Evaluation}
\label{sec:track2}

\subsection{Deployment Architecture}

\subsubsection{Python/PyTorch: NFS-Augmented Docker Strategy}

A central challenge in containerizing PyTorch inference services is the size of the ML runtime. A standard NVIDIA CUDA-enabled PyTorch container, which is what was used in this study for GPU-accelerated inference, occupies approximately 3.98--4.02~GB per service instance. While this includes the full GPU stack necessary for fast inference, it creates significant operational costs: slow registry pulls, high cold-start latency, and large storage requirements.

To address the dependency management challenge without abandoning GPU capability, an NFS-augmented architecture was developed. Rather than reinstalling Python ML libraries inside each container via \texttt{pip install}, a central NFS server VM exports a shared directory (\texttt{/external-libs/}) containing pre-built virtual environments for each model (e.g., \texttt{MNIST\_venv/}, \texttt{LSTM\_venv/}). Each Docker container mounts this shared directory at runtime and overrides \texttt{PYTHONPATH} to point to the NFS-hosted \texttt{site-packages}, as illustrated in Fig.~\ref{fig:nfs_arch}.

\begin{figure}[htbp]
\centering
\includegraphics[width=\linewidth]{figures/fig10_nfs_architecture.png}
\caption{NFS-augmented Docker deployment architecture for Python/PyTorch. A central NFS server VM exports per-model virtual environments. Each inference container mounts the shared library volume at runtime, eliminating redundant dependency installation across services.}
\label{fig:nfs_arch}
\end{figure}

The Dockerfile for this approach (CPU variant) uses \texttt{python:3.12-slim} as the base image and installs only minimal system dependencies:

\begin{lstlisting}[language=Dockerfile, caption={Python CPU inference Dockerfile using NFS-mounted libraries}, label=lst:py_dockerfile, belowskip=0pt, aboveskip=4pt]
FROM python:3.12-slim
ENV PYTHONDONTWRITEBYTECODE=1
ENV PYTHONUNBUFFERED=1
ENV PYTHONPATH=/external-libs/LSTM_env/lib/\
python3.12/site-packages
WORKDIR /app
RUN apt-get update && apt-get install -y \
    libgomp1 curl \
    && rm -rf /var/lib/apt/lists/*
COPY app.py model.py config.py dataset.py ./
COPY lstm_train_python/model.pt \
     ./generated/model.pt
EXPOSE 8000
CMD ["python", "-m", "uvicorn", "app:app", \
     "--host", "0.0.0.0", "--port", "8000"]
\end{lstlisting}

Using this slim-base approach with NFS-mounted libraries reduces the container image from $\sim$2~GB (if \texttt{torch} were pip-installed) to approximately 62--75~MB. The GPU-enabled production containers that served the inference benchmarks use the NVIDIA CUDA base image, resulting in the 3.98--4.02~GB sizes reported in Table~\ref{tab:docker_sizes}. The CPU slim approach demonstrates that the Python dependency footprint can be reduced to near-Rust levels when GPU support is not required, though it was not benchmarked in this study.

Key trade-offs of the NFS approach include:
\begin{itemize}
    \item \textbf{Pros:} Eliminates redundant installs, ensures consistent library versions across all containers, and enables instant library updates without image rebuilds.
    \item \textbf{Cons:} Introduces a runtime dependency on network connectivity to the NFS server VM; CUDA driver compatibility must still be managed at the host level; NFS latency is higher than local disk for library loading.
\end{itemize}

\subsubsection{Rust/Burn: Statically Compiled Microservice}

The Rust deployment pipeline requires no equivalent workaround because Burn compiles model weights, HTTP routing logic, and inference code into a single statically linked binary. A multi-stage Dockerfile builds the release binary inside a full Rust toolchain container, then transfers the resulting artifact into a minimal \texttt{alpine:3.23} or \texttt{nvidia/vulkan:1.3-470} runtime image. No Python interpreter, pip packages, or dynamic library resolution occurs at container startup.

HTTP routing uses the Axum framework. Model weights are loaded via Burn's \texttt{CompactRecorder} from \texttt{.mpk} files serialized during training. Tensor dimension validation occurs at compile time via Rust's generic type system (\texttt{Tensor<B,~3>} vs.\ \texttt{Tensor<B,~2>}), which eliminates a class of shape-mismatch bugs that would manifest as runtime errors in Python.

\subsection{HTTP Load Testing with Locust}

\subsubsection{Setup}

Both services were subjected to identical 5-minute Locust campaigns with 100 concurrent virtual users targeting the \texttt{POST /predict} endpoint. Tests were conducted independently for each of the four model types across both backends. Each test was repeated eight times; results reported are means with variance. Zero failures were recorded across all tests for both languages.

\subsubsection{Results}

Tables~\ref{tab:locust_python} and~\ref{tab:locust_rust} present complete request statistics and latency percentile distributions for Python and Rust services respectively.

\begin{table*}[t]
\centering
\caption{Python Inference Service: Locust Load Test Results (100 Users, 5 Min, 0 Failures)}
\label{tab:locust_python}
\begin{tabular}{|l|c|c|c|c|c|c|c|c|c|}
\hline
\textbf{Model} & \textbf{Requests} & \textbf{RPS} & \textbf{Avg (ms)} & \textbf{Min} & \textbf{Max} & \textbf{P50} & \textbf{P90} & \textbf{P95} & \textbf{P99} \\
\hline
Regression   & 175,491 & 584.66 & 15.09  & 2   & 552  & 11  & 28  & 41  & 79  \\
MNIST CNN    & 161,484 & 538.12 & 30.75  & 3   & 1793 & 27  & 54  & 65  & 93  \\
Text Class.  & 73,583  & 245.16 & 250.66 & 5   & 666  & 250 & 290 & 310 & 390 \\
LSTM         & 42,241  & 140.77 & 548.36 & 8   & 8781 & 550 & 620 & 640 & 690 \\
\hline
\end{tabular}
\end{table*}

\begin{table*}[t]
\centering
\caption{Rust Inference Service: Locust Load Test Results (100 Users, 5 Min, 0 Failures)}
\label{tab:locust_rust}
\begin{tabular}{|l|c|c|c|c|c|c|c|c|c|}
\hline
\textbf{Model} & \textbf{Requests} & \textbf{RPS} & \textbf{Avg (ms)} & \textbf{Min} & \textbf{Max} & \textbf{P50} & \textbf{P90} & \textbf{P95} & \textbf{P99} \\
\hline
Regression   & 176,356 & 587.46 & 14.51   & 2   & 206  & 11   & 26   & 35   & 77   \\
MNIST CNN    & 89,445  & 329.87 & 146.32  & 8   & 572  & 140  & 220  & 250  & 320  \\
Text Class.  & 24,794  & 82.62  & 1038.78 & 12  & 1462 & 1000 & 1200 & 1200 & 1300 \\
LSTM         & 15,026  & 50.08  & 1808.99 & 154 & 5436 & 1700 & 3000 & 3400 & 4200 \\
\hline
\end{tabular}
\end{table*}

\subsubsection{Throughput Analysis}

Fig.~\ref{fig:rps} shows the throughput comparison across all four models.

\begin{figure}[htbp]
\centering
\includegraphics[width=\linewidth]{figures/fig5_rps_comparison.png}
\caption{Inference throughput (RPS) under 100-user concurrent Locust load. Python and Rust are nearly identical for regression; Python leads significantly for compute-heavy models.}
\label{fig:rps}
\end{figure}

For the regression task, both services achieve near-identical throughput (584.66 vs.\ 587.46~RPS), confirming that for lightweight inference, HTTP overhead dominates and language choice is immaterial. For MNIST, Python achieves $1.63{\times}$ higher throughput (538.12 vs.\ 329.87~RPS). The gap is most pronounced for sequential and attention-based models: Python outperforms Rust by $2.97{\times}$ for text classification (245.16 vs.\ 82.62~RPS) and by $2.81{\times}$ for LSTM (140.77 vs.\ 50.08~RPS). These gaps reflect PyTorch's mature, highly optimized kernel library (\texttt{libtorch}, cuDNN, MKL-DNN) rather than any fundamental language-level limitation.

\subsubsection{Latency Analysis}

Fig.~\ref{fig:latency} and Fig.~\ref{fig:percentiles} present the average latency and percentile distributions respectively.

\begin{figure}[htbp]
\centering
\includegraphics[width=\linewidth]{figures/fig6_latency_comparison.png}
\caption{Average inference latency (ms, log scale). For compute-heavy models (LSTM, Text Classification), the Python advantage grows substantially.}
\label{fig:latency}
\end{figure}

\begin{figure}[htbp]
\centering
\includegraphics[width=\linewidth]{figures/fig7_percentile_dist.png}
\caption{Latency percentile distribution (P50/P95/P99) across all four model types. The tail-latency advantage of Python is most visible in the LSTM task, where P99 is 690~ms (Python) vs.\ 4200~ms (Rust)---a $6.1{\times}$ gap.}
\label{fig:percentiles}
\end{figure}

Examining tail latency (P99) reveals that Rust's variance under load is higher than Python's for complex models. For LSTM, the P99 latency gap between Rust (4200~ms) and Python (690~ms) represents a $6.1{\times}$ difference---substantially larger than the mean latency gap, suggesting that under peak concurrent load the Rust LSTM service experiences more variance. For regression, P99 latencies are nearly identical (77~ms Python vs.\ 79~ms Rust), confirming parity for lightweight models.

\FloatBarrier

\subsection{Container Image Size}

\begin{figure}[htbp]
\centering
\includegraphics[width=\linewidth]{figures/fig9_docker_sizes.png}
\caption{Docker container image sizes. Python GPU containers consistently occupy $\sim$4~GB; Rust containers occupy $\sim$1~GB, a $4{\times}$ reduction.}
\label{fig:docker_sizes}
\end{figure}

\begin{table}[htbp]
\centering
\caption{Docker Container Image Size Comparison}
\label{tab:docker_sizes}
\begin{tabular}{|l|c|c|c|}
\hline
\textbf{Task} & \textbf{Python (GB)} & \textbf{Rust (GB)} & \textbf{Ratio} \\
\hline
MNIST CNN           & 3.98  & 1.02          & $3.9{\times}$ \\
Text Classification & 4.02  & $\approx$1.00 & $4.0{\times}$ \\
Regression          & 3.98  & 0.97          & $4.1{\times}$ \\
LSTM                & 3.98  & 0.97          & $4.1{\times}$ \\
\hline
\textbf{Mean}       & \textbf{3.99} & \textbf{0.99} & $\mathbf{4.0{\times}}$ \\
\hline
\end{tabular}
\end{table}

The Python containers use the NVIDIA CUDA base image with a full GPU stack. The Rust containers use a minimal Vulkan-enabled runtime image with only the compiled binary and model weights. The $4{\times}$ size difference translates directly into reduced registry pull latency, lower storage costs, and faster Kubernetes pod startup times.

\FloatBarrier

\subsection{Model Artifact Size}

\begin{figure}[htbp]
\centering
\includegraphics[width=\linewidth]{figures/fig8_model_sizes.png}
\caption{Serialized model artifact sizes (KB, log scale). For the regression model, both formats are compact (4~KB Rust vs.\ 5~KB Python). For large models, Rust is 45--48\% smaller.}
\label{fig:model_sizes}
\end{figure}

\begin{table}[htbp]
\centering
\caption{Serialized Model Artifact Size Comparison}
\label{tab:model_sizes}
\begin{tabular}{|l|c|c|c|}
\hline
\textbf{Model} & \textbf{Rust (.mpk)} & \textbf{Python (.pt)} & \textbf{Reduction} \\
\hline
MNIST CNN   & 1.1~MB  & 2.0~MB  & 45.0\% \\
Text Class. & 21~MB   & 40.6~MB & 48.3\% \\
Regression  & 4~KB    & 5~KB    & 20.0\% \\
LSTM        & 60~KB   & 120~KB  & 50.0\% \\
\hline
\end{tabular}
\end{table}

For the regression model, both artifacts are negligibly small (4~KB Rust, 5~KB Python), as the 641-parameter model contains insufficient weight data for either format's overhead to dominate. The more operationally significant differences appear at scale: text classification serializes to 21~MB (Rust) versus 40.6~MB (Python) 48\% reduction driven by PyTorch's embedded optimizer state and Python pickle metadata. Burn's \texttt{CompactRecorder} stores only weight tensors in a flat binary encoding with no embedded metadata.

\FloatBarrier

%=======================================================================
\section{Discussion}
\label{sec:discussion}

\subsection{Training Feasibility and Convergence}

Track~1 results establish that Burn is a viable end-to-end training platform across all four evaluated architectures. Final validation accuracy on MNIST is $98.52\%$ (Rust, 95\% CI: $[98.36\%, 98.68\%]$) versus $98.20\%$ (Python, 95\% CI: $[98.04\%, 98.36\%]$). These confidence intervals do not overlap, indicating a small but statistically detectable difference in favor of Rust on this task, though the practical magnitude ($0.32\%$) is negligible.

For text classification, Python achieves a higher macro F1 ($79.03\%$, CI: $[78.23\%, 79.83\%]$) than Rust ($76.51\%$, CI: $[75.87\%, 77.15\%]$), with non-overlapping confidence intervals. This difference while modest is statistically detectable and may reflect differences in weight initialization defaults or the interaction between the Noam scheduler and Burn's optimizer implementation.

For LSTM, Rust achieves a substantially lower final validation loss ($17.59$, CI: $[14.73, 20.45]$) compared to Python ($48.92$, CI: $[43.19, 54.65]$). These intervals do not overlap, suggesting a genuine difference; however, this gap is most likely attributable to differences in synthetic data generation and normalization between the two pipelines rather than a fundamental framework advantage.

Across all tasks, both frameworks converge without numerical instability. Zero NaN events were recorded in any run, confirming that Burn's automatic differentiation and gradient accumulation are numerically stable across the tested architectures.

\subsection{The Inference Throughput Gap: Technical Explanation}

The most counter-intuitive finding of this study is that Rust, despite eliminating Python's GIL and offering language-level concurrency guarantees, achieves lower inference throughput than Python for compute-intensive models. For LSTM, Python achieves $2.81\times$ higher RPS; for text classification, $2.97\times$ higher RPS. This result requires careful interpretation.

The bottleneck in ML inference is not thread scheduling or request dispatching---it is tensor kernel throughput. PyTorch's \texttt{libtorch} backend integrates cuDNN for convolutions, MKL-DNN for CPU matrix operations, and fused attention kernels developed over years of production use. Burn's GPU backend, while architecturally sound, currently lacks: (1)~\emph{fused kernels} for multi-head attention, layer normalization, and LSTM gating, which PyTorch executes in a single GPU dispatch; (2)~\emph{cuDNN integration} providing hardware-specific convolution algorithms; and (3)~\emph{dynamic request batching} at the HTTP layer.

The tail latency disparity amplifies this picture. For the LSTM task, Python's P99 latency is 690~ms versus Rust's 4200~ms, a $6.1\times$ gap substantially wider than the mean gap ($3.3\times$). This widening is consistent with GPU dispatch queue saturation under concurrent load: when GPU kernel dispatch lacks the throughput to clear the queue at the rate imposed by 100 concurrent users, requests accumulate and tail latency grows super-linearly. For regression, both achieve $\sim$585~RPS, confirming that the performance gap is a function of kernel maturity relative to model complexity, not a general property of the language or runtime.

\subsection{Deployment Advantages: Quantified}

Rust's deployment advantages are consistent across all four models. Container images average $0.99$~GB versus $3.99$~GB for Python a $4.0\times$ reduction translating to faster registry pulls, lower cold-start latency, and reduced storage costs~\cite{vasiliev2024scaling}.

For the regression model, both artifacts are negligibly small (4~KB Rust, 5~KB Python), as the 641-parameter model contains insufficient weight data for either format's overhead to dominate. The more operationally significant differences appear at scale: text classification serializes to 21~MB (Rust) versus 40.6~MB (Python) 48\% reduction driven by PyTorch's embedded optimizer state and Python pickle metadata.

\subsection{NFS Strategy: Scope and Limitations}

The NFS-augmented deployment architecture addresses the dependency management challenge for Python inference containers. Rather than reducing the deployed image size, the GPU containers remain at 3.98--4.02~GB as the NVIDIA CUDA base image is required for accelerated inference the NFS approach eliminates redundant library installation across container instances and provides a single point of version control for ML dependencies. A CPU-only slim container using \texttt{python:3.12-slim} with NFS-mounted \texttt{torch} reduces the image to approximately 62--75~MB, at the cost of GPU acceleration.

This strategy is most appropriate for controlled, on-premises deployments where NFS infrastructure can be maintained reliably. For public cloud or serverless deployments, it introduces a network dependency that may not be acceptable. Kubernetes-based deployments with NFS persistent volume claims represent a portable middle ground~\cite{k8sml2024}.


\subsection{Limitations}

Several limitations bound the scope of the conclusions. First, all experiments use the Burn GPU backend; Burn also supports a LibTorch backend that would be expected to substantially close the inference throughput gap. Second, the load testing campaigns do not include request batching, which disproportionately benefits GPU-accelerated backends. Third, training comparisons do not include profiling traces, preventing direct attribution of training time differences to specific operations. Fourth, energy consumption and carbon footprint---increasingly relevant metrics for production ML systems~\cite{strubell2019energy}---are not measured.

\subsection{Practical Decision Framework}

Based on the empirical results, the following guidance is offered for practitioners.

\textit{Choose Rust/Burn when:} deployment footprint is the primary constraint---edge, embedded, or bandwidth-limited environments where $<$1~GB containers and kilobyte-scale model artifacts are material requirements; compile-time tensor dimension safety is required to reduce production incidents from shape mismatches; or memory-safe, dependency-minimal binaries are required for security-critical deployments.

\textit{Choose Python/PyTorch when:} maximum inference throughput is the primary requirement (demonstrated $3\times$ RPS advantage for sequential and attention-based models); lowest latency is essential (demonstrated $6\times$ P99 advantage for LSTM); or the Hugging Face and torchvision ecosystems are required for pre-trained model access.

\textit{A hybrid strategy}---training in PyTorch for ecosystem richness and iteration speed, followed by weight export and reloading in a Rust/Burn inference service for deployment efficiency---may represent the best practical trade-off as Burn's cross-framework import capabilities mature.

%=======================================================================
\section{Conclusion}
\label{sec:conclusion}

This paper presented a structured, two-track empirical evaluation of Rust (Burn) and Python (PyTorch) for machine learning across training, inference deployment, and system-level characteristics. Four model architectures were evaluated---CNN, Transformer, LSTM, and feed-forward regression---with all experiments repeated eight times to report mean and variance.

The principal findings are as follows. First, training accuracy is functionally equivalent between Burn and PyTorch across all evaluated tasks, establishing Rust as a viable end-to-end training platform. Second, Python inference services outperform Rust in throughput for compute-intensive models by up to $3.0{\times}$, and in P99 tail latency by up to $6.1{\times}$, a result driven by PyTorch's mature kernel library rather than any language-level property. Third, Rust provides a consistent $4{\times}$ reduction in container image size and 45--50\% smaller model artifacts for large models, with both regression artifacts being negligibly small at this scale. Fourth, an NFS-augmented deployment strategy for Python is described that decouples library dependencies from container images, applicable in controlled on-premises environments.

These findings suggest that the two ecosystems are best understood as complementary rather than competing. Python/PyTorch dominates for throughput-sensitive serving and training workflows requiring ecosystem breadth. Rust/Burn is the superior choice for minimal-footprint, memory-safe, and edge-optimized deployment. As the Burn framework matures its GPU-accelerated backends---particularly LibTorch and native CUDA integration---the inference performance gap documented here is expected to narrow, potentially making Rust a viable primary platform across the full ML lifecycle.

%=======================================================================
\begin{thebibliography}{00}

\bibitem{paszke2019pytorch}
A. Paszke \textit{et al.}, ``PyTorch: An imperative style, high-performance deep learning library,'' in \textit{Adv. Neural Inf. Process. Syst.}, vol.~32, 2019.

\bibitem{rust2023}
The Rust Programming Language Community, ``The Rust programming language,'' 2023. [Online]. Available: \url{https://www.rust-lang.org}

\bibitem{burn2024}
N. Gagnon-Marchand \textit{et al.}, ``Burn: A flexible and comprehensive deep learning framework in Rust,'' GitHub, 2024. [Online]. Available: \url{https://github.com/tracel-ai/burn}

\bibitem{vaswani2017attention}
A. Vaswani \textit{et al.}, ``Attention is all you need,'' in \textit{Adv. Neural Inf. Process. Syst.}, vol.~30, 2017.

\bibitem{hochreiter1997lstm}
S. Hochreiter and J. Schmidhuber, ``Long short-term memory,'' \textit{Neural Computation}, vol.~9, no.~8, pp.~1735--1780, 1997.

\bibitem{kingma2015adam}
D. P. Kingma and J. Ba, ``Adam: A method for stochastic optimization,'' in \textit{Proc. ICLR}, 2015.

\bibitem{ba2016layernorm}
J. L. Ba, J. R. Kiros, and G. E. Hinton, ``Layer normalization,'' \textit{arXiv preprint arXiv:1607.06450}, 2016.

\bibitem{merkel2014docker}
D. Merkel, ``Docker: Lightweight Linux containers for consistent development and deployment,'' \textit{Linux Journal}, vol.~2014, no.~239, 2014.

\bibitem{gujarati2020clockwork}
A. Gujarati \textit{et al.}, ``Serving DNNs like clockwork: Performance predictability from the bottom up,'' in \textit{Proc. 14th USENIX OSDI}, 2020.

\bibitem{locust}
J. Hamren \textit{et al.}, ``Locust: An open-source load testing tool,'' GitHub. [Online]. Available: \url{https://github.com/locustio/locust}

\bibitem{bezanson2017julia}
J. Bezanson, A. Edelman, S. Karpinski, and V. B. Shah, ``Julia: A fresh approach to numerical computing,'' \textit{SIAM Review}, vol.~59, no.~1, pp.~65--98, 2017.

\bibitem{onnxruntime}
Microsoft, ``ONNX Runtime: Cross-platform, high performance ML inferencing and training,'' 2021. [Online]. Available: \url{https://onnxruntime.ai}

\bibitem{pytorchvstf2025}
Z. Ba Alawi, ``A comparative survey of PyTorch vs TensorFlow for deep learning: Usability, performance, and deployment trade-offs,'' \textit{arXiv preprint arXiv:2508.04035}, 2025.

\bibitem{vella2023rust}
V. Vella, ``Boosting machine learning performance with Rust,'' \textit{Better Programming}, June 2023. [Online]. Available: \url{https://betterprogramming.pub/boosting-machine-learning-performance-with-rust-aab1f3ae1424}

\bibitem{crespo2023rust}
J. Crespo, ``Supercharging machine learning performance: 4x improvement with Rust over Python,'' \textit{LinkedIn}, May 2023.

\bibitem{athanx2024rust}
A. X. Athan, ``Choosing the right Rust machine learning framework: Candle, Burn, DFDX, or tch-rs?,'' \textit{Medium}, March 2024.

\bibitem{li2024mobileinference}
Z. Li, M. Paolieri, and L. Golubchik, ``A benchmark for ML inference latency on mobile devices,'' in \textit{Proc. 7th Int. Workshop on Edge Systems, Analytics and Networking (EdgeSys)}, ACM, 2024.

\bibitem{derosa2024modelserving}
P. De Rosa \textit{et al.}, ``On the cost of model-serving frameworks: An experimental evaluation,'' in \textit{Proc. IEEE Int. Conf. Cloud Engineering (IC2E)}, 2024, pp.~221--232.

\bibitem{buyya2009cloud}
R. Buyya, C. S. Yeo, S. Venugopal, J. Broberg, and I. Brandic, ``Cloud computing and emerging IT platforms: Vision, hype, and reality for delivering computing as the 5th utility,'' \textit{Future Generation Computer Systems}, vol.~25, no.~6, pp.~599--616, 2009.

\bibitem{vasiliev2024scaling}
G. Ramesh \textit{et al.}, ``A comprehensive review on scaling machine learning workflows using cloud technologies and DevOps,'' \textit{IEEE Access}, 2024. [Online]. Available: \url{https://ieeexplore.ieee.org/document/11126113}

\bibitem{mldaas2024}
K. P. Ravikumar, N. Ahmed, and M. S. Singh, ``ML-DaaS: An integrated ML training and deployment framework for hybrid cloud,'' \textit{IEEE Access}, 2024. [Online]. Available: \url{https://ieeexplore.ieee.org/document/11261485}

\bibitem{mohanty2024cicd}
M. V. Velumani and P. K. Muthukamatchi, ``Cloud-native software defect prediction: Leveraging machine learning models in scalable CI/CD pipelines,'' \textit{IEEE Access}, 2024. [Online]. Available: \url{https://ieeexplore.ieee.org/document/11212348}

\bibitem{k8sml2024}
D. Kutsa, ``Using machine learning for analyzing performance metrics in Kubernetes,'' \textit{IJSRED}, vol.~8, no.~2, 2025. [Online]. Available: \url{https://www.ijsred.com/volume8/issue2/IJSRED-V8I2P143.pdf}

\bibitem{strubell2019energy}
E. Strubell, A. Ganesh, and A. McCallum, ``Energy and policy considerations for deep learning in NLP,'' in \textit{Proc. ACL}, 2019.

\bibitem{watada2019containers}
J. Watada \textit{et al.}, ``Emerging trends, techniques and open issues of containerization: A review,'' \textit{IEEE Access}, vol.~7, pp.~152443--152472, 2019.
 
\end{thebibliography} 
\EOD
\end{document} 